% Generated draft: requires mathematical and source review.
\documentclass[11pt]{article}
\usepackage[margin=25mm]{geometry}
\usepackage{fontspec}
\setmainfont{Latin Modern Roman}
\usepackage{amsmath,amssymb,mathtools,unicode-math}
\setmathfont{Latin Modern Math}
\usepackage{enumitem,needspace,xurl,hyperref}
\setlength{\parindent}{0pt}
\setlength{\parskip}{5pt}
\setlength{\emergencystretch}{3em}
\newcommand{\N}{\mathbb{N}}
\newcommand{\Z}{\mathbb{Z}}
\newcommand{\Q}{\mathbb{Q}}
\newcommand{\R}{\mathbb{R}}
\newcommand{\C}{\mathbb{C}}
\newcommand{\F}{\mathbb{F}}
\newcommand{\A}{\mathbb{A}}
\newcommand{\PP}{\mathbb{P}}
\newcommand{\E}{\mathbb{E}}
\newcommand{\eps}{\varepsilon}
\newcommand{\dd}{\,\mathrm{d}}
\begin{document}
\section*{UnsolvedMath ID: 30003953}
\textbf{Source problem code:} catalogue entry
\subsection{Definitions and mathematical statement}
Let \(S_g\) be a closed connected orientable surface of genus \(g\geq 2\), and put \(\Gamma=\pi_1(S_g)\). Its Gromov boundary \(\partial_\infty\Gamma\) is homeomorphic to \(S^1\). For an integer \(d\geq 2\), write \(\mathbb P(\mathbb R^d)\) for the space of one-dimensional linear subspaces of \(\mathbb R^d\).

Here a \(1\)-Anosov representation means a projective Anosov representation \(\rho:\Gamma\to\operatorname{PGL}(d,\mathbb R)\). One standard equivalent characterization is the following uniform singular-value gap: after fixing a finite generating set of \(\Gamma\) and a Euclidean norm on \(\mathbb R^d\), there exist constants \(c>0\) and \(C\geq 0\) such that
\[
\log\frac{\sigma_1(A_\gamma)}{\sigma_2(A_\gamma)}\geq c|\gamma|-C
\qquad\text{for every }\gamma\in\Gamma.
\]
Here \(|\gamma|\) is word length, \(A_\gamma\in\operatorname{GL}(d,\mathbb R)\) is any representative of \(\rho(\gamma)\), and \(\sigma_1(A_\gamma)\geq\sigma_2(A_\gamma)\) are its two largest singular values. Their ratio is independent of the scalar ambiguity in \(A_\gamma\). Such a representation has a canonical continuous equivariant boundary embedding
\[
\xi_\rho:\partial_\infty\Gamma\longrightarrow\mathbb P(\mathbb R^d).
\]
Its circular limit set is \(\Lambda_\rho=\xi_\rho(\partial_\infty\Gamma)\).

Fix any smooth Riemannian metric on projective space. The circle \(\Lambda_\rho\) is called rectifiable when its one-dimensional Hausdorff measure is finite. This property is independent of the chosen smooth metric. For \(0<\alpha\leq 1\), call \(\Lambda_\rho\) a \(C^{1,\alpha}\) embedded circle if it is the image of a \(C^1\) embedding \(f:S^1\to\mathbb P(\mathbb R^d)\) with nowhere-vanishing derivative whose first derivative is locally \(\alpha\)-Hölder in smooth coordinates. Thus failure of this regularity concerns the embedded set, rather than merely a particular boundary parametrization.

\textbf{Existence question, with unresolved parameter scope.} For parameters \(g,d,\alpha\) as above, consider the assertion
\[
\exists\,\rho:\pi_1(S_g)\to\operatorname{PGL}(d,\mathbb R)
\quad\text{such that}\quad
\begin{cases}
\rho\text{ is }1\text{-Anosov},\\
\mathcal H^1(\Lambda_\rho)<\infty,\\
\Lambda_\rho\text{ is not a }C^{1,\alpha}\text{ embedded circle}.
\end{cases}
\]
The supplied problem asks for a representation of this kind. It does not specify whether \(d\) is prescribed or may be chosen, nor whether \(\alpha\) is prescribed or the same limit circle must fail to be \(C^{1,\alpha}\) for every positive Hölder exponent. The displayed assertion is therefore a parameterized formulation, not a resolution of those missing quantifiers. The range \(0<\alpha\leq1\) is the usual Hölder convention used here; the source supplies no range. The wording also leaves implicit whether existence is requested for every genus or for some genus. No irreducibility assumption is imposed in the supplied question.
\subsection{Short English statement}
The question asks whether a projective Anosov representation of a closed orientable surface group can have a limit circle of finite length that lacks \(C^{1,\alpha}\) regularity. The dimension, the scope of the genus quantifier, and the quantification over Hölder exponents require clarification. The supplied record labels the problem open; this draft does not independently verify that status.
\subsection{Sources}
\begin{itemize}
\item Supplied local registry record in \texttt{data/unsolvedmath/problems.json}, permanent UnsolvedMath ID \(30003953\), title \emph{Rectifiable Nonsmooth Limit Curves of Anosov Representations}, exact source problem code \texttt{catalogue entry}; collection: Oberwolfach Reports Open Problems.
\item \emph{New Trends in Teichmüller Theory and Mapping Class Groups} (2019), Oberwolfach Reports, pp.\ 2475--2534, DOI \texttt{10.4171/owr/2018/40}. Supplied literature-source label: ``Oberwolfach Report 40, Zhang problem.'' \url{https://doi.org/10.4171/owr/2018/40}. The record describes the report as documenting workshop open-problem discussions. Its proposer field is empty; the Zhang attribution is retained only as supplied in the reference label.
\item Zhang--Zimmer, supplied reference label ``regularity of Anosov limit sets'': \url{https://arxiv.org/abs/1903.11021}.
\item Literature assessment dated exactly \texttt{2026-08-21}: supplied dated curation, not independently verified. The record labels the problem open and states that Zhang--Zimmer give necessary and sufficient regularity criteria for irreducible surface-group representations, but that their work does not settle the requested existence of a rectifiable limit circle lacking \(C^{1,\alpha}\) regularity. The supplied curation explicitly requests independent verification of time-sensitive status and cited claims.
\end{itemize}
\end{document}
